% AUTHOR: Amlal El Mahrouss
% PURPOSE: Glossary of Definitions — merged from parts 1–11.

\documentclass[10pt]{article}

%% ---- Core Packages ----------------------------------------
\usepackage[T1]{fontenc}
\usepackage[utf8]{inputenc}
\usepackage{lmodern}
\usepackage{microtype}

%% ---- Page Geometry ----------------------------------------
\usepackage[
  top=1in, bottom=1in,
  left=0.75in, right=0.75in,
  columnsep=0.25in
]{geometry}

%% ---- Mathematics ------------------------------------------
\usepackage{amsmath, amssymb, amsthm, amsfonts}
\usepackage{mathtools}
\usepackage{bm}
\usepackage{textcomp}
\usepackage{esint}

%% ---- Figures & Tables -------------------------------------
\usepackage{graphicx}
\usepackage{booktabs}
\usepackage{multirow}
\usepackage{array}
\usepackage{caption}
\usepackage{subcaption}
\usepackage{float}

%% ---- Code Listings ----------------------------------------
\usepackage{listings}
\usepackage{xcolor}
\usepackage{algorithm}
\usepackage{algpseudocode}

\lstdefinestyle{codestyle}{
  backgroundcolor=\color{gray!8},
  basicstyle=\ttfamily\footnotesize,
  breakatwhitespace=false,
  breaklines=true,
  captionpos=b,
  commentstyle=\color{green!50!black},
  keywordstyle=\color{blue}\bfseries,
  stringstyle=\color{orange!70!black},
  numberstyle=\tiny\color{gray},
  numbers=left,
  numbersep=5pt,
  showstringspaces=false,
  frame=single,
  rulecolor=\color{gray!40},
  tabsize=2
}
\lstset{style=codestyle}

%% ---- Hyperlinks -------------------------------------------
\usepackage[
  colorlinks=true,
  linkcolor=blue!70!black,
  citecolor=green!50!black,
  urlcolor=blue!60!black
]{hyperref}
\usepackage{url}

%% ---- Bibliography -----------------------------------------
\usepackage[numbers, sort&compress]{natbib}

%% ---- Miscellaneous ----------------------------------------
\usepackage{enumitem}
\usepackage{siunitx}
\usepackage{cleveref}

%% ---- Theorem Environments ---------------------------------
\theoremstyle{plain}
\newtheorem{theorem}{Theorem}[section]
\newtheorem{lemma}[theorem]{Lemma}
\newtheorem{corollary}[theorem]{Corollary}
\newtheorem{proposition}[theorem]{Proposition}

\theoremstyle{definition}
\newtheorem{definition}[theorem]{Definition}
\newtheorem{example}[theorem]{Example}

\theoremstyle{remark}
\newtheorem{remark}[theorem]{Remark}

%% ---- Custom Commands --------------------------------------
\newcommand{\R}{\mathbb{R}}
\newcommand{\N}{\mathbb{N}}
\newcommand{\Z}{\mathbb{Z}}
\newcommand{\E}{\mathbb{E}}
\newcommand{\Prob}{\mathbb{P}}
\newcommand{\norm}[1]{\left\lVert #1 \right\rVert}
\newcommand{\abs}[1]{\left\lvert #1 \right\rvert}
\newcommand{\inner}[2]{\left\langle #1,\, #2 \right\rangle}
\newcommand{\ie}{\textit{i.e.}\xspace}
\newcommand{\eg}{\textit{e.g.}\xspace}
\newcommand{\etal}{\textit{et al.}\xspace}

%% ---- Caption Formatting -----------------------------------
\captionsetup{
  font=small,
  labelfont=bf,
  format=hang,
  justification=justified
}

\usepackage{lettrine}

%% ---- Header / Footer --------------------------------------
\usepackage{fancyhdr}
\pagestyle{fancy}
\fancyhf{}
\fancyhead[R]{\small\thepage}
\renewcommand{\headrulewidth}{0.4pt}

%% ============================================================
%%  FRONT MATTER
%% ============================================================

\title{%
	\vspace{-1.5em}%
	\large\textbf{Glossarium Aequationum.}%
}
\author{By A. E. Mahrouss}
\date{\today}

%% ============================================================
\begin{document}
%% ============================================================

\maketitle
\tableofcontents
\newpage

%% ============================================================
\section{Introduction:}
%% ============================================================

\lettrine[lines=2, lhang=0.1, loversize=0.2]{T}he document's purpose is to provide a reference for theorems, definitions, and equations for the purpose of being reusability. References are cited at the bottom of the glossary where credit is due.\\
There will be proofs part of another book, adjacent to this one. Where each steps will be explained.\\
One caveat: $y(x), \lambda(x), \rho(x)$ are defined generically in this book. However the Lean proofs will have a concrete function attributed to them.

%% ============================================================
\section{Definitions:}
%% ============================================================

\lettrine[lines=1, lhang=0.1, loversize=0.2]{L}et following theorems, definitions, and remarks, be defined in $\mathbb{R}.$\\The following are defined:

\begin{theorem}
Let the following:
$\forall S \in \mathbb{R}, \quad Baud(MIN) < Freq(S) < Baud(MAX).$ Where $Baud(MIN)$ is the minimal baud rate and $Baud(MAX)$ the maximal baud rate:
\begin{align*}
	Freq(S) &= \frac{Hi(S)}{Lo(S)} Weight(x),\\
	-\frac{Freq(S)}{Weight(S)} &= \frac{Hi(S)}{Lo(S)}.
\end{align*}
By Hartley's Law, the weight of $S$ over its frequency $F$, may be equal of the division between the highest signal and lowest signal of $S$.
\end{theorem}

\begin{theorem}
	Let for any $\forall a, \mu, \lambda, b, c \in \mathbb{R}$:
	\begin{align*}
		V &= \int (\lambda + \mu)\sqrt{ax^2 + bx  +c} \operatorname{dx}.
	\end{align*}
	We define three general equations for $\forall V'$:
	\begin{align*}
		V' &= (\lambda + \mu)\sqrt{ax^2 + bx  +c}x \int x \operatorname{d}(\lambda + \mu)\sqrt{ax^2 + bx  +c},\\
		\frac{V'}{(\lambda + \mu)} &= \sqrt{ax^2 + bx  +c}x \int x \operatorname{d}(\lambda + \mu)\sqrt{ax^2 + bx  +c},\\
		\sqrt{ax^2 + bx  +c}\frac{V'}{(\lambda + \mu)} &= (\sqrt{ax^2 + bx  +c})^{2}x \int x \operatorname{d}(\lambda + \mu)\sqrt{ax^2 + bx  +c}.
	\end{align*}
\end{theorem}

\begin{theorem}
	Thus let the following per Theorem 2.2:
	\begin{align*}
		O' + (\lambda + \mu + x)^{n-2} &= (\lambda + \mu + x)^{n-1} + (\lambda + \mu + x)^{n},\\
		\frac{O'}{(\lambda + \mu + x)^{n-2}} &= \frac{(\lambda + \mu + x)^{n-1} + (\lambda + \mu + x)^{n}}{(\lambda + \mu + x)^{n-2}}, \quad O' \neq 0.
	\end{align*}
	Let $\forall O \in \mathbb{R}$ be:
	\begin{align*}
		O &= \frac{d}{dx} O', \quad \forall O, O' \in \mathbb{R}, \quad O' \neq 0.
	\end{align*}
\end{theorem}

\begin{theorem}
	Let the following for any $\forall m, n \in \mathbb{R}^{+}.$:
\begin{equation}
    \mu(n, p) \coloneq \sum_{n}^{p} M(n, p), \quad \mu(n, p) \in \mathbb{D}^{+}
\end{equation}
	$\forall \mu(n, p)$ There exists:
\begin{equation}
    \mu(n, p) = p! + \text{ } (n+1p)! + \text{ } (n+2p)! + \text{ } \dots \text{ } + (n+xp)!.
\end{equation}
\end{theorem}

\begin{theorem}
	Let a set of rational numbers $\forall x \in \mathbb{R},$ in $\forall \mu(n, p) \in \mathbb{R},$ be in the following inequality:
	\begin{align*}
		\frac{2\pi}{\mu(m, p)} < x < \mu(n, p), \quad x \neq 0.
	\end{align*}
\end{theorem}

\begin{theorem}
	Let for any real $\forall x \in \mathbb{R}.$ If: $x \neq 0.$ Then:
	\begin{align*}
		E(x) &= \frac{\gamma(x)}{\delta(x) + \lambda(x)}, \quad E(x) \in \mathbb{R}, \quad \gamma(x) \neq 0.
	\end{align*}
	Therefore, based on the definition $E(x)$ in $\mathbb{R}.$:
	\begin{align*}
		P &= \prod_{n=1}^{3}\rho(x) x^{n} + \int E(x) dx,\\
		P &= \prod_{n=1}^{3}\rho(x) x^{n} + \int \frac{\gamma(x)}{\delta(x) + \lambda(x)} dx,\\
		P &= \int \frac{\gamma(x)}{\delta(x) + \lambda(x)} dx + x^{n} \rho(x)^{3}, \quad P(x) \neq 0.
	\end{align*}
	Separately, let the following for $\forall x, n \in \mathbb{R}$:
	\begin{align*}
		\prod_{n=1}^{3}\rho(x) x^{n} + \int \frac{\gamma(x)}{\delta(x) + \lambda(x)} dx \in \mathbb{R}.
	\end{align*}
	There exists a solution in $\forall s = P' \in \mathbb{R},$ of the following:
	\begin{align*}
		P' &= \prod_{n=1}^{3}\rho(x) x^{n} + \int \frac{\gamma(x)}{\delta(x) + \lambda(x)} dx, \quad P' \neq 0.
	\end{align*}
	Separately, let the following $\forall x, L \in \mathbb{R}.$:
	\begin{align*}
		L + \lambda(x)^{2} &= \frac{\gamma(x)}{\delta(x)},\\
	 	L + \lambda(x) &= \frac{\gamma(x)}{\delta(x)  \lambda(x) },\\
	 	L &= \frac{\gamma(x)}{\delta(x)  \lambda(x) } - \lambda(x), \quad \delta(x)\lambda(x) \neq 0.
	\end{align*}
	Where the value of $\forall \Delta y \in \mathbb{R}$ equals:
	\begin{align*}
		&= \frac{\ln x}{\Delta y},\\
		\Delta y &= \ln x,
	\end{align*}
	Separately, there exists for $\forall I, r, x, \Delta x, n \in \mathbb{R}$:
	\begin{align*}
		I &\coloneq -r + ( \frac{\ln x}{\Delta n} + \frac{\Delta n}{\ln x})^m, \quad r \neq 0, \quad r \in \mathbb{R}^{+}.
	\end{align*}
	For all variables of $\forall r \in \mathbb{R}.$ Being equal to each other, e.g: $x = n = \ln x = m.$ Then the following equation:
	\begin{align*}
		r &= ( \frac{\ln x}{n} + \frac{n}{\ln x})^m, \quad r \neq 0, \quad r \in \mathbb{R}{+}.
	\end{align*}
	Admits a transcendental number, If and only If: $r \neq 0, \quad r \in \mathbb{R}$.
\end{theorem}

\begin{theorem}
	Let the following plots both a real and imiginary part for both. If: $x, n \in \mathbb{R}.$ Then:
	\begin{align*}
		\frac{7}{\log n} + \frac{363 \log 7}{140} &= \sum_{n=1}^{m} \frac{\ln x}{n} + \frac{n}{\ln x}, \quad m = 7, \quad \log n \neq \log 1.
	\end{align*}
	And thus by computing the products of $(\frac{\ln x}{n} + \frac{n}{\ln x}).$:
	\begin{align*}
		\frac{7}{\log n} + \frac{\log^{7} 7}{5040} &= \prod_{n=1}^{m} \frac{\ln x}{n} + \frac{n}{\ln x}, \quad m = 7, \quad \log n \neq \log 1.
	\end{align*}
\end{theorem}

\begin{theorem}
	Let us define the following equations for $\forall n \in \mathbb{R}^{+}$:
	\begin{align*}
		Z &= \frac{\phi(n^{s}!)}{\tau(n^{s}!)} + \frac{\phi(n^{s}!)}{\lambda(n^{s}!)} + \phi(n^{s}!), \quad Z \in \mathbb{Z}.
	\end{align*}
	That admits the following alternate form:
	\begin{align*}
		Z' &= \frac{\phi(n^{s}!)}{\tau(n^{s}!)} + \frac{\phi(n^{s}!)}{\lambda(n^{s}!)} + \phi(n^{s}!),\\
		&= \phi(n^{s}!)(\frac{1}{\tau(n^{s}!)} + \frac{1}{\lambda(n^{s}!)}) + \phi(n^{s}!),\\
		&= \phi(n^{s}!)(\frac{1}{\tau(n^{s}!)} + \frac{1}{\lambda(n^{s}!)} + 1).
	\end{align*}
	Alternatively, we can assume another equation:
	\begin{align*}
		\phi(s, n) &= \frac{(s-3)}{\tau(n^{s}!)} + \frac{(s-2)}{\lambda(n^{s}!)} + \frac{(s-1)}{f'(n^{s}!)} + s, \quad s \neq 0.
	\end{align*}
	Assuming $\phi(s, n) \neq 0.$ And $\phi(s, n) \in \mathbb{R}.$
\end{theorem}

\begin{definition}
	Let a theorem of the following equations for any $\forall C, x, n \in \mathbb{R}$:
	\begin{align*}
		C &= \int \frac{\lambda(y) \sqrt{x}}{x} dx = - \prod_{n=1}^{y}\rho(x) x^{n} + \mathrm{C}',\\
		C &= \int \frac{\lambda(y) \sqrt{x}}{x} dx + \prod_{n=1}^{y}\rho(x) x^{n} = \mathrm{C}'
	\end{align*}
	Admits a set of results in a domain $\mathbb{D}^{+}$. Alternatively we can suppose the following equations $\pm E \in \mathbb{R}.$
	\begin{align*}
		\gamma(x)+\lambda(x) + \prod_{n=2}^{4}\frac{1}{\rho(x)}^{n} &= \pm \mathrm{E},\\
		\gamma(x)+\lambda(x) + \frac{1}{\rho(x)} \frac{1}{\rho(x)}^{2} \frac{1}{\rho(x)}^{3} \frac{1}{\rho(x)}^{4} &= \pm \mathrm{E},\\
		\gamma(x)+\lambda(x) + \frac{1}{\rho(x)^{10}} &= \pm \mathrm{E}.
	\end{align*}
	If, assuming all variables are non-zero.
\end{definition}

\begin{definition}
	Let us define the following for $\forall \Delta x, x, y \in \mathbb{R}$:
	\begin{align*}
	\Delta x &=	y + \gamma(x), \quad \gamma(x) < \Delta x.
	\end{align*}
  	Separately, we have the following for $x \in \mathbb{R}$:
	\begin{align*}
		\mathrm{A} &= \gamma(x)+\Delta x + \frac{1}{\lambda(x)}, \quad \lambda(x) \neq 0.
	\end{align*}
  Separately, we have the following  for $\forall x, y \in \mathbb{R}$:
	\begin{align*}
		\frac{x}{\lambda(x)} + \gamma(x) &= \frac{x}{y} + \mathrm{Y'},\\
		\frac{x}{\lambda(x)} + \gamma(x) - \frac{x}{y} &= \mathrm{Y'}, \quad \Delta x \in \mathbb{D},\\
		\lambda(x) + \gamma(x) -\frac{x}{y} &= \frac{\mathrm{Y'}}{x}, \quad x \neq 0.
	\end{align*}
	For all valid numbers in a domain '$\mathbb{D}$'.
\end{definition}

\begin{theorem}
	Let an equation of the following for $\forall L, x, y, z \in \mathbb{R}$:
	\begin{align*}
		L &= \iiint (x y z)^{3} dx.dy.dz,\\
		c_{1} y z + c_{2} z + c_{3} + \frac{1}{64} + x^{4} + y^{4} + z^{4} &= \iiint (x y z)^{3} dx.dy.dz.
	\end{align*}
	Admits the following equation $\forall c, x, y, z \in \mathbb{R}.$:
	\begin{align*}
		(c_{1} y z + c_{2} z + c_{3} + \frac{1}{64} + x^{4} + y^{4} + z^{4})^{n-2} + \sum_{n=1}^{p} \frac{ (c_{1} y z + c_{2} z + c_{3} + \frac{1}{64} + x^{4} + y^{4} + z^{4})^{n}}{ (c_{1} y z + c_{2} z + c_{3} + \frac{1}{64} + x^{4} + y^{4} + z^{4})^{n-1}} &= \pm \mathrm{C},
	\end{align*}
\end{theorem}

\begin{definition}
	We define Euler–Mascheroni constant:
	\begin{equation}
		\gamma = \lim_{n \to \infty}(\sum_{k=1}^{n}(\frac{1}{k}-\log(n))).
	\end{equation}
	Which may be rewritten to per M. Ram Murty and N. Saradha findings:
	\begin{equation}
		\gamma(a, q) = \lim_{n \to \infty}\sum_{k=0}^{n}(\frac{1}{a+kq}-\frac{\log(a+nq)}{q}).
	\end{equation}
\end{definition}

\begin{definition}
	The following is the closed form of a series of low-integer-order polylogarithms:
	\begin{align*}
		\frac{\operatorname{d}}{\operatorname{dz}} \operatorname{Li}_{n}(z) &= \frac{\operatorname{Li}_{n-1}(z)}{z}.
	\end{align*}
\end{definition}

\begin{definition}
	The following is a series of the exponential function:
	\begin{align*}
		\sum_{k=0}^{\infty} \frac{z^k}{k!} &= e^{z}.
	\end{align*}
\end{definition}

\begin{theorem}
	The following is a series of Euler–Mascheroni constants over $k$:
	\begin{align*}
		\sum_{k=1}^{\infty} \frac{k}{\gamma(a, qk)}, \quad qk \neq 0.
	\end{align*}
\end{theorem}

\begin{definition}
	The following is a series of $k$ over Euler–Mascheroni constants:
	\begin{align*}
		\sum_{k=1}^{\infty} \frac{\gamma(a, qk)}{k}, \quad k \neq 0.
	\end{align*}
\end{definition}

\begin{definition}
	The following inequality may also be defined:
	\begin{align*}
		\sum_{k=1}^{\infty} \frac{\gamma(a, qk)}{k} < x < \sum_{k=1}^{\infty} \frac{k}{\gamma(a, qk)}, \quad x \neq 0.
	\end{align*}
\end{definition}

\begin{definition}
	Applying the previous theorem laid out on previous papers. We satisfy the following:
	\begin{align*}
		y &= \pm z \frac{b}{c!}, \quad \pm z \in \mathbb{D}^{+}\\
		\frac{y}{\pm z} &= c!, \quad \pm z \neq 0\\
		-c! &= \pm \frac{y}{\pm z}.
	\end{align*}
	For any domain $\mathbb{D}^{+}$. Where the following is a Taylor series of:
	\begin{equation}
		z = \frac{x^n}{n!}, \quad n \neq 0.
	\end{equation}
\end{definition}

\begin{theorem}
	Let for $\forall y, \pm z, a, q, k, x \in \mathbb{R}$:
	\begin{align*}
		y &= \sum_{k=1}^{\infty} \frac{\gamma(a, qk)}{k} < x < \sum_{k=1}^{\infty} \frac{k}{\gamma(a, qk)},\\
		\frac{y}{\pm z} &= \frac{\sum_{k=1}^{\infty} \frac{\gamma(a, qk)}{k} < x < \sum_{k=1}^{\infty} \frac{k}{\gamma(a, qk)}}{\pm z}.
	\end{align*}
	Now in order to continue further, we need to solve $\forall x \in \mathbb{R}.$:
	\begin{align*}
		\sum_{k=1}^{\infty} \frac{\gamma(a, qk)}{k} &\leq x < \sum_{k=1}^{\infty} \frac{k}{\gamma(a, qk)},\\
		\sum_{k=1}^{\infty} \frac{\gamma(a, qk)}{k} &\leq x < \sum_{k=1}^{\infty} \frac{k}{\gamma(a, qk)},\\
		\sum_{k=1}^{\infty} \frac{\gamma(a, qk)}{k} + \sum_{k=1}^{\infty} \frac{k}{\gamma(a, qk)} &\leq x,
	\end{align*}
	Therefore, for $x = \sum_{k=1}^{\infty} \frac{\gamma(a, qk)}{k} + \sum_{k=1}^{\infty} \frac{k}{\gamma(a, qk)}, \forall x \in \mathbb{R}.$
	\begin{align*}
		\pm z + y &= \sum_{k=1}^{\infty} \frac{\gamma(a, qk)}{k} + \sum_{k=1}^{\infty} \frac{k}{\gamma(a, qk)},\\
		y &= \frac{\sum_{k=1}^{\infty} \frac{\gamma(a, qk)}{k} + \sum_{k=1}^{\infty} \frac{k}{\gamma(a, qk)}}{\pm z}.
	\end{align*}
	Then for: x = $\sum_{k=1}^{\infty} \frac{\gamma(a, qk)}{k} + \sum_{k=1}^{\infty} \frac{k}{\gamma(a, qk)}, \forall x \in \mathbb{R}.$:
	\begin{align*}
		\pm z + \frac{\sum_{k=1}^{\infty} \frac{\gamma(a, qk)}{k} + \sum_{k=1}^{\infty}\frac{k}{\gamma(a, qk)}}{\pm z} &= \sum_{k=1}^{\infty} \frac{\gamma(a, qk)}{k} + \sum_{k=1}^{\infty}\frac{k}{\gamma(a, qk)},\\
		\frac{\sum_{k=1}^{\infty} \frac{\gamma(a, qk)}{k} + \sum_{k=1}^{\infty} \frac{k}{\gamma(a, qk)}}{\pm z} &= \frac{\sum_{k=1}^{\infty} \frac{\gamma(a, qk)}{k} + \sum_{k=1}^{\infty} \frac{k}{\gamma(a, qk)}}{\pm z}, \quad \pm z \neq 0.
	\end{align*}
\end{theorem}

\begin{definition}
	We define the Golden-ratio constant:
	\begin{align*}
		\phi &= \frac{1+\sqrt{5}}{2}, \forall \phi \in \mathbb{R}.
	\end{align*}
	And the Supergolden-ratio constant:
	\begin{align*}
		\psi &= \frac{1+\frac{29+3\sqrt{93}}{2}+\frac{29+3\sqrt{93}}{2}}{3}, \forall \psi \in \mathbb{R}.
	\end{align*}
	And the Kepler-Bouwkamp constant:
	\begin{align*}
		K &= \prod_{n=3}^{\pi} \cos(\frac{\pi}{n}),\\
		&= \cos(\frac{\pi}{n-1}) \cos(\frac{\pi}{n}) \dots, \forall K \in \mathbb{R}.
	\end{align*}
	And. For an equation of the previous result per Theorem 2.19:
	\begin{align*}
		A &= \pm z + \frac{\sum_{k=1}^{\infty} \frac{\gamma(a, qk)}{k} + \sum_{k=1}^{\infty}\frac{k}{\gamma(a, qk)}}{\pm z},\\
		\pm z^2 &= \sum_{k=1}^{\infty} \frac{\gamma(a, qk)}{k} + \sum_{k=1}^{\infty}\frac{k}{\gamma(a, qk)},\\
		\mathrm{A}' = \pm z^2 &= \sum_{k=1}^{\infty} \frac{\gamma(a, qk)}{k} + \sum_{k=1}^{\infty}\frac{k}{\gamma(a, qk)},\\
		\mathrm{A}'' = \pm z^{2} - \sum_{k=1}^{\infty}\frac{k}{\gamma(a, qk)} &= \sum_{k=1}^{\infty} \frac{\gamma(a, qk)}{k}, \quad k \neq 0, \quad \forall A \in \mathbb{R}.
	\end{align*}
\end{definition}

\begin{theorem}
	Let for the following result for $\forall x \in \mathbb{R}.$:
	\begin{align*}
		\sqrt{2}+2&=\sqrt{2} + x,\\
		2\sqrt{2}+2 &=x,\\
		2(\sqrt{2} + 1) &= x,
	\end{align*}
	Then we can solve for $l=x, \forall l \in \mathbb{R}.$
	\begin{align*}
		l + r &\leq n - 1,\\
		2(\sqrt{2} + 1) + r &\leq n - 1.
	\end{align*}
	Also. $\forall x$ There exists for $r \in \mathbb{D}^{+}.$
	\begin{align*}
		2(\sqrt{2} + 1) + \frac{\frac{n - 1}{2} - 4}{2}  &\leq \frac{n - 1}{2} + n - 1,\\
		2(\sqrt{2} + 1) + \frac{\frac{n - 1}{2} - 4}{2} - \frac{n - 1}{2} &\leq n - 1.
	\end{align*}
\end{theorem}

\begin{theorem}
	Let a general theorem of the following for $\forall x \in \mathbb{R}$ be:
	\begin{align*}
		\mathrm{B} &= 2(\sqrt{2} + 1) + \frac{\frac{n - 1}{2} - 4}{2} - \frac{n - 1}{2},\\
		&= \frac{\frac{\frac{n - 1}{2} - 4}{2} - \frac{n - 1}{2}}{2(\sqrt{2} + 1)}.
	\end{align*}
	Then. Substituting $\frac{n - 1}{2}=\phi, \quad \forall \phi \in \mathbb{R}.$:
	\begin{align*}
		\mathrm{C} &= 2(\sqrt{2} + 1) + \frac{\frac{n - 1}{2} - 4}{2} - \phi,\\
		&= \frac{\frac{\phi - 4}{2} - \phi}{2(\sqrt{2} + 1)},\\
		\mathrm{C}' &= \frac{\phi - 4}{2} - \phi.
	\end{align*}
\end{theorem}

\begin{theorem}
	Let the following constant $C \in \mathbb{R}.$ admit a derivative of the following:
	\begin{align*}
		\mathrm{C}' &= \frac{\phi - 4}{2} - \phi,\\
		\mathrm{C}' &= \frac{\frac{1+\sqrt{5}}{2} - 4}{2} - \frac{1+\sqrt{5}}{2}.
	\end{align*}
\end{theorem}

\begin{theorem}
	Let the following equation for $\forall n \in \mathbb{R}.$
	\begin{align*}
		2(\sqrt{2} + 1) + \frac{\frac{n - 1}{2} - 4}{2}  &\leq \frac{n - 1}{2} + n - 1,\\
		2(\sqrt{2} + 1) + \frac{\frac{n - 1}{2} - 4}{2} - \frac{n - 1}{2} &\leq n - 1,\\
	\end{align*}
	 That admits the following solution of $(2(\sqrt{2} + 1) + \frac{\frac{n - 1}{2} - 4}{2}) \in \mathbb{R}.$:
	\begin{align*}
	 	2(\sqrt{2} + 1) + \frac{\frac{n - 1}{2} - 4}{2} &\leq 0.
	\end{align*}
\end{theorem}

\begin{definition}
	Let the following equation for $\forall n \in \mathbb{R}.$:
	\begin{align*}
		2(\sqrt{2} + 1) + \frac{\frac{n - 1}{2} - 4}{2} &= x,
	\end{align*}
	Thus for the following variable $n=5, \quad \forall n \in \mathbb{N}^{+}.$:
	\begin{align*}
		2(\sqrt{2} + 1) + \frac{5}{2} &= x,\\
		\sqrt{2}+2 + \frac{5}{2} &= x.
	\end{align*}
	And the following for $x \in \mathbb{R}.$:
	\begin{align*}
		2+\frac{5}{2} &= \frac{x}{\sqrt{2}}.
	\end{align*}
\end{definition}

\begin{theorem}
	Let the following derivative, and primitive function for $\forall x \in \mathbb{R}.$:
	\begin{align*}
		f'(x)&= \frac{\ln (x)^{\mathrm{e}}}{\mathrm{e}} \cdot \mathrm{e}^{x},\\
		f(x) &= \int_{0}^{1} f'(x)dx,\\
		&= \int_{0}^{1} \frac{\ln (x)^{\mathrm{e}}}{\mathrm{e}} \cdot \mathrm{e}^{x} dx.
	\end{align*}
\end{theorem}

\begin{theorem}
	Let the following equation for $\forall x \in \mathbb{R}.$:
	\begin{align*}
		x &= 2(\sqrt{2} + 1) + \frac{5}{2}.
	\end{align*}
	Thus for $\forall x$:
	\begin{align*}
		\int_{0}^{1} \frac{\ln (x)^{\mathrm{e}}}{\mathrm{e}} \cdot \mathrm{e}^{x} dx.
	\end{align*}
	Then there exists the following definite integral admits a solution $\mathrm{R} \in \mathbb{R}.$
	\begin{align*}
		\mathrm{R} &= \int_{0}^{1} \frac{\ln (x)^{\mathrm{e}}}{\mathrm{e}} \cdot \mathrm{e}^{x} dx,\\
		&= \int_{0}^{1} \frac{\ln (2(\sqrt{2} + 1) + \frac{5}{2})^{\mathrm{e}}}{\mathrm{e}} \cdot \mathrm{e}^{2(\sqrt{2} + 1) + \frac{5}{2}},\\
		&= \mathrm{e}^{7/2+2 \sqrt{2}} \log^{e}(\frac{9}{2}+2+\sqrt{2}) + \mathrm{C}.
	\end{align*}
\end{theorem}

\begin{theorem}
	Let the following definite integral for $\forall x, y, \pm D \in \mathbb{R}.$:
	\begin{align*}
		\int_{C}^{\pi} \frac{ \ln |x| \cdot f(y)f'(x) }{x^{n}y!} dx &= \pm D, \quad \pm D \neq 0.
	\end{align*}
	With a definite integral that admits a solution of:
	\begin{align*}
		\int_{C}^{\pi} \frac{ \ln |x| \cdot f(y)f'(x) }{x^{n}y!} dx &= [\frac{ \ln |x| \cdot f(y)f'(x) }{x^{n}y!}]_{C}^{\pi} + D.
	\end{align*}
	Let us impose two inequalities for $\forall x, y, n \in \mathbb{R}.$:
	\begin{align*}
		\frac{\pi}{\frac{ \ln |x| \cdot f(y)f'(x) }{x^{n}y!}} < \frac{ \ln |x| \cdot f(y)f'(x) }{x^{n}y!} < \sum_{n=1}^{\infty} \frac{ \ln |x| \cdot f(y)f'(x) }{x^{n}y!},\\
	\frac{\pi}{\ln |x| \cdot f(y)f'(x) } < \ln |x| \cdot f(y)f'(x)  < \sum_{n=1}^{\infty} \ln |x| \cdot f(y)f'(x).
	\end{align*}
\end{theorem}

\begin{theorem}
	Let the following equation for $\forall x, y, n, C \in \mathbb{R}.$:
	\begin{align*}
		\pm \frac{\gamma(x-y)}{|x - y|^{n}}! &= \pm \frac{\sum_{n=1}^{p}(\frac{\gamma(x-y)}{1})}{\sum_{n=1}^{p}(|x - y|^{n})} + C, \quad x \neq y, \quad y \neq 0, \quad x \neq 0.
	\end{align*}
\end{theorem}

\begin{definition}
	Let the transform for $N \neq 0, \quad N \in \mathbb{N}.$ Be the following:
	\begin{align*}
		X_k &= \sum_{m=0}^{N-1} x_{m}\mathrm{e}^{-i2\pi \frac{km}{N}n}, \quad k = 0,\dots,N-1.
	\end{align*}
	The inverse notation for $n \neq 0, N \neq 0, \quad N \in \mathbb{N}.$ Is given by:
	\begin{align*}
		x_k &= \frac{1}{N} \sum_{m=0}^{N-1} X_k \mathrm{e}^{i2\pi \frac{km}{N}n}, \quad k = 0,\dots,N-1.
	\end{align*}
	For $x_{0}, \dots, x_{n-1}$ be complex numbers. Where $\mathrm{e}^{-i2\pi /n}$ is a primitive nth root of 1.
\end{definition}

\begin{theorem}
	Let the following identitiy in $\mathbb{R}.$:
	\begin{align*}
		\frac{ln(\pi^\mathrm{e})!}{\pi} &= (\mathrm{e} \log(\pi))!.
	\end{align*}
	And for $\mathrm{H} \in \mathbb{R}.$:
	\begin{align*}
		\mathrm{H} &\coloneq \frac{ln(\pi^\mathrm{e})!}{\pi} + \frac{(\mathrm{e} \log(\pi))!}{\pi}, \quad \forall \mathrm{H} \in \mathbb{R}.
	\end{align*}
	From the following for $x = \frac{ln(\pi^\mathrm{e})!}{\pi}, \quad \forall x \in \mathbb{R}.$:
	\begin{align*}
	\mathrm{H} &\coloneq \frac{ln(\pi^\mathrm{e})!}{\pi} + \frac{(\mathrm{e} \log(\pi))!}{\pi}, \quad \forall \mathrm{H} \in \mathbb{R}.
	\end{align*}
	We then would establish another identity $\forall x \in \mathbb{R}.$:
	\begin{align*}
		x &= -\frac{(\log(\pi)!) (e \log(\pi))!}{\pi}, \quad x \neq 0.
	\end{align*}
\end{theorem}

\begin{definition}
	Defining the Chvátal–Sankoff constants for $\gamma_k, n, k \in \mathbb{R}.$:
	\begin{equation}
		\gamma_k = \lim_{n \to \infty} \frac{\mathrm{E}{\lambda_{n, k}}}{n}, \quad n \neq 0.
	\end{equation}
\end{definition}

\section{Final Remarks:}
\lettrine[lines=2, lhang=0.1, loversize=0.2]{T}he previous section, which establishes 31 definitions and theorems. Will now serve as a basis for plotting.\\
Which in the next chapter 'On Theorems' will be used to plot some of our theorems. But before, let us impose the following remark in $\mathbb{R}.$:

\begin{remark}
	Per Theorem 2.31, for $x = \frac{ln(\pi^\mathrm{e})!}{\pi}, \quad \forall x \in \mathbb{R}.$ We have the following:
	\begin{equation}
		\frac{ln(\pi^\mathrm{e})!}{\pi} = \frac{(\mathrm{e} \log(\pi))!}{\pi}, x \neq 0, \quad \forall x \in \mathbb{R}.
	\end{equation}
\end{remark}

\section{Conclusion:}

\lettrine[lines=2, lhang=0.1, loversize=0.2]{T}he nine pages of theorems now comes to an end, stating them here was the first step to then develop a framework later in the series of papers. We will next plot and analyse our results based of our theorems and definitions. We will now proceed to plot and analyze the following theorems to analyze their convergence and general behavior.

\nocite{*}

\bibliographystyle{acm}
\bibliography{Bell_System_Hartley, S0022314X10001836, Fatou1906, A_Course_of_Pure_Mathematics, General_Theory_Dirichlet_Series, Riemann, bonito2026, hardy1922theory, chambadal1981dictionnaire}

\end{document}
